%%%%%%%%%%%%%%%%%%%%%%%%%%%%%%%%%%%%%%%%%%%%%%%%%%%%%%%%%%%%%%%%%%%%%%
% Lớp 10 | Phương pháp tọa độ trong mặt phẳng (𝑂𝑥𝑦) | Ba đường conic trong mặt phẳng tọa độ
% Dạng 0H9?5-1: Xác định các yếu tố của elip
% Số câu: 23  (sinh tự động bởi tach_cau.py; từ nay thầy sửa trực tiếp file này)
%%%%%%%%%%%%%%%%%%%%%%%%%%%%%%%%%%%%%%%%%%%%%%%%%%%%%%%%%%%%%%%%%%%%%%

%%%%%%%%%%%%%%%% TRẮC NGHIỆM 4 PHƯƠNG ÁN %%%%%%%%%%%%%%%%

% nguồn: MaHoaNH_TracNghiem_Lop10
\begin{ex}%[0H9N5-1]%[ma:f6d64df5]
	Cho elip $(E)$ có phương  trình chính tắc là
	$\dfrac{x^2}{16}+\dfrac{y^2}{4}=1$. Gọi $2c$ là tiêu  cự của
	$(E)$.  Mệnh đề nào sau đây đúng?
	\choice{\True $c^2=12$}{$c^2=16$}{$c^2=20$}{$c^2=4$} 
	\loigiai{ 
		Ta có  $a^2=16, b^2=4$. Vậy $c^2=a^2-b^2=12$.
	} 
\end{ex}

% nguồn: MaHoaNH_TracNghiem_Lop10
\begin{ex}%[0H9N5-1]%[ma:b37016a4]
	Cho  elip $(E)$ có phương trình chính tắc là $\dfrac{x^2}{5}+\dfrac{y^2}{1}=1$. Tìm toạ độ các đỉnh $A_1$, $A_2$ của $(E)$.
	\choice 
		{ $A_1(-1;0)$ và $A_2(1;0)$} 
		{$A_1(-2;0)$ và  $A_2(2;0)$}
		 {\True $A_1\left(-\sqrt{5};0\right)$ và $A_2\left(\sqrt{5};0\right)$}
		{$A_1(-4;0)$  và $A_2(4;0)$}
	\loigiai{Ta có $a=\sqrt{5}$,  do đó $A_1(-\sqrt{5};0)$ và $A_2(\sqrt{5};0)$.
	 }
\end{ex}

% nguồn: MaHoaNH_TracNghiem_Lop10
\begin{ex}%[0H9N5-1]%[ma:7ca30efe]
	Cho elip $(E)$ có phương trình chính tắc là $\dfrac{x^2}{4}+\dfrac{y^2}{1}=1$. Tìm tọa độ các đỉnh $B_1$, $B_2$  của $(E)$.
	 \choice
		{\True $B_1(-1;0)$  và $B_2(1;0)$}
		 {$B_1(1;0)$ và $B_2(-1;0)$}
		{ $B_1(-\sqrt{5};0)$  và $B_2(\sqrt{5};0)$}
		{$B_1(-2;0)$ và $B_2(2;0)$} 
	\loigiai{Ta có $b=1$,  do đó $B_1(-1;0)$ và $B_2(1;0)$.
	} 
\end{ex}

% nguồn: MaHoaNH_TracNghiem_Lop10
\begin{ex}%[0H9N5-1]%[ma:d2826742]
	Cho elip $(E)$ có phương trình  chính tắc là $\dfrac{x^2}{25}+\dfrac{y^2}{16}=1$. Tìm độ dài trục bé của $(E)$.
	 \choice
		 { $10$}
		 {\True $8$}
		 {$9$}
		 {$6$}
	\loigiai{Ta  có $b=4$, do đó độ dài trục bé của $(E)$ là $2b=8$.
	} 
\end{ex}

% nguồn: MaHoaNH_TracNghiem_Lop10
\begin{ex}%[0H9N5-1]%[ma:e1b02579]
	Cho elip $(E)$ có phương trình chính tắc là $\dfrac{x^2}{5}+\dfrac{y^2}{1}=1$. Tìm tọa độ các tiêu điểm của  $(E)$.
	\choice 
		{ $F_1(-1;0)$  và $F_2(1;0)$}
		{\True $F_1(-2;0)$ và $F_2(2;0)$} 
		{$F_1(2;0)$  và $F_2(-2;0)$}
		{$F_1(\sqrt{5};0)$ và  $F_2(-\sqrt{5};0)$}
	\loigiai{Ta có $a=\sqrt{5}$, $b=1$, suy ra  $c=2$ do đó $F_1(-2;0)$ và $F_2(2;0)$.
	 }
\end{ex}

% nguồn: MaHoaNH_TracNghiem_Lop10
\begin{ex}%[0H9N5-1]%[ma:00eae849]
	Cho elip $(E)$  có phương trình chính tắc là $\dfrac{x^2}{4}+\dfrac{y^2}{1}=1$. Gọi $F_1$, $F_2$ là các tiêu điểm của $(E)$ và điểm $M$ tùy ý trên $(E)$. Mệnh đề nào đúng?
	\choice 
		{  $MF_1+MF_2=\sqrt{2}$}
		{$MF_1+MF_2=6$} 
		{\True $MF_1+MF_2=4$} 
		 {$MF_1+MF_2=1$}
	\loigiai{Ta có $a=2$, do đó $MF_1+MF_2=2a=4$. 
	 }
\end{ex}

% nguồn: MaHoaNH_TracNghiem_Lop10
\begin{ex}%[0H9N5-1]%[ma:a73705ef]
	Cho elip $(E)$ có phương trình chính tắc là $\dfrac{x^2}{3^2}+\dfrac{y^2}{2^2}=1$. Tìm  độ dài trục lớn của $(E)$.
	\choice 
		{$ 4  $}
		{\True  $ 6 $}
		{$ 5 $} 
		{$ 9 $} 
	 \loigiai{
	Ta có $ a=3 $, $ b=2 $ suy ra độ dài trục lớn $  2a=6 $.
	} 
\end{ex}

% nguồn: MaHoaNH_TracNghiem_Lop10
\begin{ex}%[0H9N5-1]%[ma:ea6fb393]
	Trong mặt phẳng $Oxy$, cho elip có phương trình chính tắc $\dfrac{x^2}{36}+\dfrac{y^2}{16}=1$. Tiêu  cự của elip đã cho bằng
	 \choice
		{$8$} 
		{$12$} 
		 {$2\sqrt{5}$}
		 {\True $4\sqrt{5}$}
	 \loigiai{ Từ $\dfrac{x^2}{36}+\dfrac{y^2}{16}=1$ suy ra $a^2=36$, $b^2=16$.\\
		Ta có  $c^2=a^2-b^2=36-16=20\Rightarrow c=\sqrt{20}=2\sqrt{5}$.\\
		 Tiêu cự của elip đã cho là $F_1F_2=2c=4\sqrt{5}$.
	} 
\end{ex}

% nguồn: MaHoaNH_TracNghiem_Lop10
\begin{ex}%[0H9N5-1]%[ma:22e8556d]
	Cho elip có phương trình $\dfrac{x^2}{25}+\dfrac{y^2}{16}=1$. Elip có tiêu cự  bằng
	 \choice
		 {$16$}
		{$4$} 
		{\True $6$} 
		 {$12$}
	\loigiai{ 
		Ta có $a=5$, $b=4$,  $c=\sqrt{a^2-b^2}=\sqrt{5^2-4^2}=3$.\\
		 Tiêu cự của elip là $2c=6$.
	} 
\end{ex}

% nguồn: MaHoaNH_TracNghiem_Lop10
\begin{ex}%[0H9N5-1]%[ma:df7c78e3]
	Cho elip $(E)\colon \dfrac{x^2}{9}+\dfrac{y^2}{4}=1$. Tìm  khẳng định \textbf{sai} trong các khẳng định sau.
	\choice 
		{$(E)$ có tiêu điểm $F_1\left(\sqrt{5}; 0\right),  F_2\left(\sqrt{5}; 0\right)$}
		{$(E)$ có tỉ số $\dfrac{c}{a}=\dfrac{\sqrt{5}}{3}$} 
		 {$(E)$ có đỉnh $A_1(-3; 0)$}
		{\True $(E)$ có độ dài  trục lớn là $ 3 $}
	 \loigiai{
	Ta có $ a=3 $, $ b=2 $. Do  đó $ c=\sqrt{9-4}=\sqrt{5} $.\\
	Do đó $(E)$ có  tiêu điểm $F_1\left(-\sqrt{5}; 0\right)$, $F_2\left(\sqrt{5}; 0\right)$.\\
	Tỉ số $  \dfrac{c}{a}=\dfrac{\sqrt{5}}{3} $.\\
	$(E)$ có đỉnh  $A_1(-3; 0)$.\\
	$(E)$ có độ dài trục lớn  là $ 2a=6 $.
	 }
\end{ex}

% nguồn: MaHoaNH_TracNghiem_Lop10
\begin{ex}%[0H9N5-1]%[ma:564374cb]
	Cho elip $(E)\colon x^2+4 y^2=1$. Tìm mệnh đề 
	đúng trong các  mệnh đề sau.
	\choice 
		{$(E)$ có độ dài trục lớn bằng $ 4 $} 
		{$ (E) $ có độ dài trục bé bằng $ 2 $} 
		{\True $(E)$ có đỉnh  $A_1(-1; 0)$}
		{$(E)$  có tiêu cự bằng $2\sqrt{3}$}
	 \loigiai{
	Ta viết lại $ (E)\colon \dfrac{x^2}{1}+\dfrac{y^2}{\frac{1}{4}} =1$. 
	Khi đó $ a=1 $, $ b=\dfrac{1}{2} $, $ c=\dfrac{\sqrt{3}}{2}  $.
	\begin{itemize}
		\item $ (E) $ có độ dài trục lớn bằng $  2 $.
		\item $ (E) $ có độ dài trục bé  bằng $ 1 $.
		\item $ (E) $ có  đỉnh là $ (-1;0) $ và $ (1;0) $.
		\item$ (E) $ có tiêu cự $ \sqrt{3} $. 
	\end{itemize}
	 }
\end{ex}

% nguồn: MaHoaNH_TracNghiem_Lop10
\begin{ex}%[0H9H5-1]%[ma:924a9a9c]
	Cho elip $(E)$ có phương trình chính tắc là $\dfrac{x^2}{25}+\dfrac{y^2}{16}=1$. Tính tính diện tích hình chữ nhật đi  qua 4 đỉnh của $(E)$.
	 \choice
		{  $20$}
		{\True $80$} 
		 {$48$}
		{$60$} 
	\loigiai{Ta có $a=5$, $b=4$, do đó diện  tích hình chữ nhật đi qua 4 đỉnh của $(E)$ là $2b\cdot 2a=80$.
	 }
\end{ex}

% nguồn: MaHoaNH_TracNghiem_Lop10
\begin{ex}%[0H9H5-1]%[ma:12fbcf67]
	Cho  elip $\left(E\right)$ có phương trình chính tắc là $\dfrac{x^2}{9}+\dfrac{y^2}{4}=1$. Gọi $A_1$, $A_2$ là các đỉnh của $(E)$ thuộc trục $Ox$. Tính độ dài đoạn thẳng $A_1A_2$.
	\choice 
		{\True  $A_1A_2=6$}
		{$A_1A_2=4$} 
		 {$A_1A_2=5$}
		 {$A_1A_2=2\sqrt{5}$}
	 \loigiai{
		Vì $(E)\colon\dfrac{x^2}{9}+\dfrac{y^2}{4}=1$ suy ra $A_1(-3;0)$ và $A_2(3;0)$. Mặt  khác $A_1A_2=2a=6$.
	 }
\end{ex}

% nguồn: MaHoaNH_TracNghiem_Lop10
\begin{ex}%[0H9H5-1]%[ma:f7ed6250]
	Cho elip $\left(E\right)$ có phương trình chính tắc là $\dfrac{x^2}{25}+\dfrac{y^2}{16}=1$. Tính tổng độ dài các trục của $(E)$. 
	 \choice
		{\True $10$} 
		{$9$} 
		 {\True $18$}
		{$8$} 
	\loigiai{ 
		Từ phương trình của elip ta suy ra trục lớn là $2a=10$, trục bé $2b=8$. Tổng độ  dài hai trục là $18$.
	} 
\end{ex}

% nguồn: MaHoaNH_TracNghiem_Lop10
\begin{ex}%[0H9H5-1]%[ma:681195ae]
	Cho elip $(E)$ có phương trình chính tắc $\dfrac{x^2}{9}+\dfrac{y^2}{4}=1$. Tổng khoảng cách từ một điểm $M$ bất kỳ trên $(E)$ đến hai tiêu điểm  là
	\choice 
		 {\True $6$}
		{$4$} 
		 {$3$}
		{$9$} 
	\loigiai{ 
		Phương trình $(E)$ có dạng $ \dfrac{x^2}{a^2}+\dfrac{y^2}{b^2}=1$, với $a,\, b>0$, $a=3$. \\ 
		Theo định nghĩa $M\in (E)\Leftrightarrow  MF_1+MF_2=2a=6$.\\
		Tổng  khoảng cách từ một điểm $M$ bất kỳ trên $(E)$ đến hai tiêu điểm bằng $6$.
	} 
\end{ex}

% nguồn: MaHoaNH_TracNghiem_Lop10
\begin{ex}%[0H9H5-1]%[ma:d6dcaa02]
	Trong mặt phẳng $Oxy$, cho Elip có phương trình  chính tắc $\dfrac{x^2}{34}+\dfrac{y^2}{9}=1$. Điểm nào dưới đây là một tiêu điểm của Elip đã cho?
	\choice 
		{$F_1\left(-3;0\right)$} 
		{\True  $F_2\left(-5;0\right)$}
		 {$F_3\left(0;-5\right)$}
		 {$F_4\left(\sqrt{34};0\right)$}
	 \loigiai{
	Từ phương  trình chính tắc của Elip đã cho ta có
	\[ 
	 \heva{&a^2=34\\ &b^2=9} \Rightarrow c=\sqrt{a^2-b^2}=\sqrt{34-9}=5.
	\] 
	Vậy $F_2\left(-5;0\right)$ là một tiêu điểm của Elip đã  cho.
	 }
\end{ex}


%%%%%%%%%%%%%%%% ĐÚNG/SAI %%%%%%%%%%%%%%%%

% nguồn: MaHoaNH_DungSai_Lop10
\begin{ex}%[0H9N5-1]%[ma:00ef9eb3]
	Cho elip  có phương trình chính tắc $\dfrac{x^2}{25}+\dfrac{y^2}{16}=1$. Trong các mệnh đề sau, mệnh đề nào đúng, mệnh đề nào sai?
	\choiceTF 
		{\True Elip có tiêu điểm là $F_1(-3;0)$} 
		{\True Elip có hai tiêu điểm là  $F_2(3;0)$}
		{\True Tiêu cự của elip  là $F_1F_2=6$}
		{\True Tổng các khoảng cách từ mỗi điểm trên elip tới hai  tiêu điểm bằng $10$}
	\loigiai{Ta có $a^2=25$, $b^2=16$.  Do đó, $c=\sqrt{a^2-b^2}=3$.\\
		Vậy $F_1(-3;0)$,  $F_2(3;0)$ và tiêu cự $F_1F_2=2c=6$.\\
		Từ đó, suy ra tổng các khoảng cách từ mỗi điểm trên elip tới hai tiêu điểm  bằng $2a=10$.
	 }
\end{ex}

% nguồn: MaHoaNH_DungSai_Lop10
\begin{ex}%[0H9N5-1]%[ma:bdc73fbc]
	Cho elip có phương trình chính tắc $\dfrac{x^2}{36}+\dfrac{y^2}{9}=1$. Trong các mệnh đề sau, mệnh đề  nào đúng, mệnh đề nào sai?
	\choiceTF 
		{\True Elip có tiêu  điểm là $F_1(-3\sqrt{3};0)$}
		{Elip có  tiêu điểm là $F_2(0;3\sqrt{3})$}
		{\True Tiêu cự của elip là $F_1F_2=6\sqrt{3}$} 
		 {Tổng các khoảng cách từ mỗi điểm trên elip tới hai tiêu điểm bằng $6$}
	\loigiai{Ta  có $a^2=36$, $b^2=9$. Do đó, $c=\sqrt{a^2-b^2}=3\sqrt{3}$.\\
		Vậy $F_1(-3\sqrt{3};0)$, $F_2(3\sqrt{3};0)$  và tiêu cự $F_1F_2=2c=6\sqrt{3}$.\\
		Từ đó, suy ra tổng các khoảng cách từ  mỗi điểm trên elip tới hai tiêu điểm bằng $2a=12$.
	 }
\end{ex}

% nguồn: MaHoaNH_DungSai_Lop10
\begin{ex}%[0H9N5-1]%[ma:808c19a1]
	Trong mặt phẳng  $Oxy$ cho Elip $(E)\colon\dfrac{x^2}{16}+\dfrac{y^2}{9}=1$.
	\choiceTF 
		{\True Độ dài trục lớn của $(E)$ bằng  $8$}
		{Độ dài trục bé của  $(E)$ bằng $3$}
		 {Tiêu cự của $(E)$ bằng $7$}
		{\True $B_1(0;3)$ là một  đỉnh của Elip của $(E)$}
	 \loigiai{
		 Với $a^2=16, b^2=9 \Rightarrow c^2=a^2-b^2=7$.\\
		Do  đó, $a=4, b=3, c=\sqrt{7}$. Từ đó suy ra:
		\begin{itemchoice}
			\itemch Độ dài trục lớn  bằng $8$;
			\itemch Độ dài trục bé bằng  $6$;
			\itemch Tiêu cự bằng $2\sqrt{7}$; 
			\itemch  $B_1(0;3)$ là một đỉnh của Elip.
		\end{itemchoice}
	} 
\end{ex}

% nguồn: MaHoaNH_DungSai_Lop10
\begin{ex}%[0H9H5-1]%[ma:d81acc9d]
	 Trong mặt phẳng $Oxy$, cho Elip $(E)$ có phương trình $\dfrac{x^2}{16}+\dfrac{y^2}{4}=1$.
	 \choiceTF
		{\True Độ dài trục lớn của $(E)$ là  $8$}
		{Tiêu cự của $(E)$ là  $12$}
		{\True  Điểm $M\left(2;\sqrt{3}\right)\in (E)$}
		{Diện tích hình chữ nhật cơ sở  của $(E)$ là $8$}
	\loigiai{ 
		\begin{itemchoice}
			\itemch Đúng. Vì độ  dài trục lớn là $2a=2\cdot 4=8$.
			\itemch Sai. Vì  $c^2=a^2-b^2=4^2-2^2=12\Rightarrow c=2\sqrt{3}$, nên độ dài trục tiêu cự là $2c=2\cdot2\sqrt{3}=4\sqrt{3}$.
			\itemch  Đúng. Thay $M\left(2;\sqrt{3}\right)$ vào phương trình $(E)$ ta được $\dfrac{2^2}{16}+\dfrac{\sqrt{3}^2}{4}=1$, nên $M\left(2;\sqrt{3}\right)\in (E)$.
			\itemch Sai. Diện  tích hình chữ nhật cơ sở là $2a\cdot2b=8\cdot4=32$.
		\end{itemchoice}
	} 
\end{ex}

% nguồn: MaHoaNH_DungSai_Lop10
\begin{ex}%[0H9H5-1]%[ma:481a938d]
	Cho elip $(E)$ có hai tiêu điểm là $F_1$, $F_2$; độ dài  trục lớn bằng $20$ và tiêu cự bằng $10$.
	\choiceTF 
		{\True Nếu $M$ là  một điểm tùy ý trên elip $(E)$ thì $MF_1+MF_2=20$}
		{\True Tọa  độ các tiêu điểm là $F_1(-5;0)$ và $F_2(5;0)$}
		{Elip $(E)$  có độ dài trục bé là $150$}
		{\True Phương  trình chính tắc của elip $(E)$ là $\dfrac{x^2}{100}+\dfrac{y^2}{75}=1$}
	 \loigiai{
		\begin{itemchoice}
			\itemch Đúng. Do độ dài trục lớn là $2a=20$ nên theo định nghĩa ta có $MF_1+MF_2=2a=20$. 
			\itemch Đúng. Do  tiêu cự $2c=10$ nên $c=5$. Suy ra, elip $(E)$ có hai tiêu điểm là $F_1(-5;0)$ và $F_2(5;0)$.
			 \itemch Sai. Ta có $c^2=a^2-b^2\Leftrightarrow b^2=a^2-c^2=10^2-5^2=75$. Suy ra, độ dài trục bé bằng $2b=10\sqrt{3}$.
			\itemch  Đúng. Phương trình chính tắc của elip cần tìm là
			$$ (E)\colon  \dfrac{x^2}{100}+\dfrac{y^2}{75}=1.$$
		\end{itemchoice}
	 }
\end{ex}

% nguồn: MaHoaNH_DungSai_Lop10
\begin{ex}%[0H9V5-1]%[ma:d43ecd6b]
	Trong mặt phẳng $Oxy$, cho Elip $(E)\colon 4x^2+16y^2=1$.  \choiceTF
	 {$(E)$ có độ dài trục lớn bằng $2$}
	{\True Điểm $K\left(0;-\dfrac{1}{4}\right)$ là một đỉnh  của $(E)$}
	{$(E)$ có tiêu cự bằng  $\sqrt{3}$}
	{\True Các tiêu điểm của $(E)$ là  $F_1\left(-\dfrac{\sqrt{3}}{4};0\right)$ và $F_2\left(\dfrac{\sqrt{3}}{4};0\right)$}
	 \loigiai{
		 Ta có $4x^2+16y^2=1\Leftrightarrow \dfrac{x^2}{\dfrac{1}{4}}+\dfrac{y^2}{\dfrac{1}{16}}=1$.\\
		Khi đó $\heva{&a^2=\dfrac{1}{4}\\&b^2=\dfrac{1}{16}}\Rightarrow  \heva{&a=\dfrac{1}{2}\\&b=\dfrac{1}{4}.}$\\
		 $b^2+c^2=a^2\Leftrightarrow c^2=a^2-b^2=\dfrac{1}{4}-\dfrac{1}{16}=\dfrac{3}{16}\Rightarrow c=\dfrac{\sqrt{3}}{4}$.
		\begin{itemchoice}
			\itemch $(E)$ có độ dài trục lớn bằng  $2a=2 \cdot \dfrac{1}{2}=1$.
			\itemch Vì $b=\dfrac{1}{4}$ nên điểm $K\left(0;-\dfrac{1}{4}\right)$ là một đỉnh của $(E)$. 
			 \itemch $(E)$ có tiêu cự bằng $2c=2\cdot \dfrac{\sqrt{3}}{4}=\dfrac{\sqrt{3}}{2}$.
			\itemch Các tiêu điểm của $(E)$ là  $F_1\left(-\dfrac{\sqrt{3}}{4};0\right)$ và $F_2\left(\dfrac{\sqrt{3}}{4};0\right)$.
		\end{itemchoice}
	} 
\end{ex}

% nguồn: MaHoaNH_DungSai_Lop10
\begin{ex}%[0H9V5-1]%[ma:dd983343]
	 Trong mặt phẳng tọa độ $Oxy$ cho elip $(E)\colon\dfrac{x^2}{25}+\dfrac{y^2}{16}=1$.
	 \choiceTF
		{Độ dài trục lớn của elip  là $5$}
		{\True Hai tiêu  điểm của elip là $F_1(-3;0)$, $F_2(3;0)$}
		{Giao điểm của elip với trục tung là các  điểm $(-5;0)$ và $(5;0)$}
		{Có đúng $8$ điểm có tọa độ  là các số nguyên thuộc $(E)$}
	 \loigiai{
		\begin{itemchoice}
			\itemch Ta có $\heva{&a=\sqrt{25}=5 \\&b=\sqrt{16}=4}\Rightarrow  c=\sqrt{a^2-b^2}=\sqrt{25-16}=3$.\\
			Độ dài  trục lớn của elip là $2a=2\cdot5=10$.
			 \itemch Elip có hai tiêu điểm là $F_1(-3;0)$, $F_2(3;0)$.
			 \itemch Giao điểm của elip với trục tung là các điểm $B_1(0;-4)$ và $B_2(0;4)$.
			 \itemch Ta có $\heva{&\dfrac{x^2}{25}\geq0 \\&\dfrac{y^2}{16}\geq0\\&\dfrac{x^2}{25}+\dfrac{y^2}{16}=1}$ nên $\dfrac{y^2}{16}\leq1\Leftrightarrow y^2\leq16$ mà $y\in\mathbb{Z}$ nên $y\in\{-4;-3;-2;-1;0;1;2;3;4\}$.\\
			Vì  $\dfrac{x^2}{25}+\dfrac{y^2}{16}=1\Rightarrow x=\pm5\sqrt{1-\dfrac{y^2}{16}}$.\\
			Chỉ có các  giá trị $y\in\{-4;0;4\}$ cho ta các giá trị $x$ nguyên.
			\begin{itemize}
				\item  Với $y=-4\Rightarrow x=0$. Ta có $1$ điểm có tọa độ nguyên thuộc $(E)$.
				\item Với $y=0\Rightarrow x=\pm5$. Ta  có $2$ điểm có tọa độ nguyên thuộc $(E)$.
				\item Với  $y=4\Rightarrow x=0$. Ta có $1$ điểm có tọa độ nguyên thuộc $(E)$.
			\end{itemize}
			Vậy có tất cả $5$  điểm có tọa độ nguyên thuộc $(E)$.
		\end{itemchoice}
	 }
\end{ex}
